\documentclass{ximera}

\begin{document}
	\author{Stitz-Zeager}
	\xmtitle{Exercises for Polar Conics}{}

\mfpicnumber{1} \opengraphsfile{ExercisesforPolarConics} % mfpic settings added 

\begin{question}
Graph the following equations.

\begin{problem}
$x^2+2xy+y^2 -x\sqrt{2}+y\sqrt{2} -6= 0$
\end{problem}
 
\begin{problem}
$7x^2-4xy\sqrt{3}+3y^2-2x-2y\sqrt{3}-5= 0$
\end{problem}

\begin{problem}
$5x^2+6xy+5y^2 - 4\sqrt{2}x+4\sqrt{2}y = 0$
\end{problem}

\begin{problem}
$x^2+ 2\sqrt{3}xy+3y^2+ 2\sqrt{3}x-2y-16 = 0$
\end{problem}

\begin{problem}
$13x^2-34xy\sqrt{3}+47y^2 - 64=0$
\end{problem}

\begin{problem}
$x^2-2\sqrt{3} xy-y^2+8=0$
\end{problem}

\begin{problem}
$x^2-4xy+4y^2-2x\sqrt{5}-y\sqrt{5}=0$
\end{problem}

\begin{problem}
 $8x^2+12xy+17y^2 - 20 = 0$
\end{problem}
\end{question}


\begin{question}
Graph the following equations.

\begin{problem}
$r = \frac{2}{1-\cos(\theta)}$
\end{problem}

\begin{problem}
 $r = \frac{3}{2 + \sin(\theta)}$
\end{problem}

\begin{problem}
$r = \frac{3}{2-\cos(\theta)}$
\end{problem}

\begin{problem}
$r = \frac{2}{1 + \sin(\theta)}$
\end{problem}

\begin{problem}
$r = \frac{4}{1+3\cos(\theta)}$
\end{problem}

\begin{problem}
$r = \frac{2}{1-2\sin(\theta)}$
\end{problem}

\begin{problem}
$r = \frac{2}{1 + \sin(\theta - \frac{\pi}{3})}$
\end{problem}

\begin{problem}
$r = \frac{6}{3 - \cos\left(\theta + \frac{\pi}{4}\right)}$
\end{problem}

\end{question}

\begin{question}
 The matrix $A(\theta) = \left[ \begin{array}{rr} \cos(\theta) & -\sin(\theta) \\ \sin(\theta) & \cos(\theta) \\ \end{array} \right]$  is called a \textbf{rotation matrix}\index{matrix ! rotation}\index{rotation matrix}. 
 
 \smallskip
 
 We've seen this matrix most recently  used in the proof of Theorem \ref{rotatecoordinatesthm}.

\begin{problem}
  Show the matrix from Example \ref{rotationmatrixex} in Section \ref{MatArithmetic} is none other than $A\left(\frac{\pi}{4}\right)$.
\end{problem}

\begin{problem}
Discuss with your classmates how to use $A(\theta)$ to rotate points in the plane.
\end{problem}

\begin{problem}
Using the even / odd identities for cosine and sine,  show $A(\theta)^{-1} = A(-\theta)$.  Interpret this geometrically.
\end{problem}

\end{question}

\end{document}